% !TeX spellcheck = it_IT

\documentclass[format=169, handout]{beamer}

\usepackage[utf8]{inputenc}
\usepackage[italian]{babel}
\usepackage{multicol}

\usetheme{metropolis}

\usepackage{mathtools}
\DeclarePairedDelimiter{\ceil}{\lceil}{\rceil}

\title{Esercitazioni di Informatica B}
\subtitle{Introduzione a Matlab}

\author{Stefano Cereda\\
	stefano.cereda@polimi.it
}
\date{13/11/2018}
\institute[PoliMi]{\vspace{0.5cm}\centering Politecnico di Milano \\ \vspace{0.2cm}
	\includegraphics[width=0.2\textwidth]{../logopolimi}}

\setbeamercovered{invisible}

\makeindex

\begin{document}
	\begin{frame}
	\maketitle
\end{frame}

\begin{frame}<handout:0>{Conclusione esercizi struct}
Completare gli esercizi sulle struct in C (anagrafe)
\end{frame}


\begin{frame}{Costrutto if: equazioni di secondo grado}
Scaricare e installare Matlab/Octave

Script: commenti(nome e descrizione; data e autori; def variabili), input, calcolo, output.

Si chieda l’inserimento dei tre coefficienti di una equazione di secondo grado: $ax^2+bx+c=0$ e si stampi il valore delle radici reali.
\end{frame}

\begin{frame}{Grafici: equazioni di secondo grado}
Si chieda l’inserimento dei tre coefficienti di una equazione di secondo grado: $ax^2+bx+c=0$ e si dia una rappresentazione della curva.
\end{frame}

\begin{frame}{Costrutto while: fattoriale}
Si scriva uno script che chieda l'inserimento di un numero $n$ e ne calcoli e visualizzi il fattoriale utilizzando un ciclo while.
\end{frame}

\begin{frame}{Matrice triangolare}
Scrivere uno script in Matlab/Octave che svolga le seguenti operazioni:
\begin{enumerate}
	\item Definire una matrice quadrata di interi di dimensione $N\times N$, con $N=5$.
	\item Acquisire da tastiera una sequenza di valori interi da memorizzare nella porzione triangolare superiore della matrice. La parte triangolare superiore della matrice è costituita dalle celle poste sulla diagonale principale e al di sopra di tale diagonale.
	\item Copiare in un array di dimensioni opportune tutti i valori della porzione triangolare di matrice che sono maggiori di -10.
	\item Stampare il contenuto dell'array.
\end{enumerate}
\end{frame}

\end{document}